\documentclass[11pt,a4paper]{article}
\usepackage[T1]{fontenc}
\usepackage{lmodern}
\usepackage[margin=27mm]{geometry}
\usepackage{amsmath,amssymb,amsthm,mathtools}
\usepackage{microtype}
\usepackage{booktabs,array}
\usepackage{enumitem}
\usepackage[hidelinks]{hyperref}
\setlength{\parskip}{5pt plus 1pt}
\setlength{\parindent}{0pt}
\setlength{\emergencystretch}{3em}
\setlist{itemsep=2pt,topsep=3pt}
\raggedbottom
\newtheorem{theorem}{Theorem}[section]
\newtheorem{lemma}[theorem]{Lemma}
\newtheorem{corollary}[theorem]{Corollary}
\newtheorem{proposition}[theorem]{Proposition}
\theoremstyle{definition}
\newtheorem{definition}[theorem]{Definition}
\newtheorem{example}[theorem]{Example}
\theoremstyle{remark}
\newtheorem{remark}[theorem]{Remark}
\newcommand{\Z}{\mathbb{Z}}
\newcommand{\N}{\mathbb{N}}
\newcommand{\Q}{\mathbb{Q}}
\newcommand{\F}{\mathbb{F}}
\newcommand{\ps}[1]{#1[[x]]}
\newcommand{\coeff}[2]{[x^{#1}]\,#2}
\newcommand{\hl}[1]{\textbf{#1}}
\newcommand{\para}[1]{\par\medskip\noindent\textbf{#1}\enspace}
\hypersetup{pdftitle={Residues modulo 16 for a series defined by a bilateral theta-type equation},pdfauthor={Dev D. Goyal},pdfsubject={OEIS A355872}}
\title{Residues modulo $16$ for a series defined by a bilateral theta-type equation\\[6pt]\large A proof of three conjectures in OEIS A355872}
\author{Dev D. Goyal}
\date{October 2026}
\begin{document}
\maketitle
\begin{abstract}
The sequence A355872 of the On-Line Encyclopedia of Integer Sequences (OEIS), contributed by Paul~D.~Hanna, lists the
coefficients $a(n)$ of the power series $A(x)=\sum_{n\ge1}a(n)x^{4n-3}=2x+14x^5+434x^9+\cdots$ determined by the
equation $x=\sum_{j\in\Z}(-x)^{j^2}A(x)^{(j-1)^2}$, a sum over all integers $j$. The entry conjectures that every $a(n)$ is
$2$ modulo $4$, and that $a(n)$ is $2$ or $6$ modulo $8$ according as $n$ is odd or even. We prove the stronger statement
$a(n)\equiv2(-1)^{n-1}\pmod{16}$ for every $n\ge1$. The proof first shows that the equation has exactly one solution over
every commutative ring and that only the degrees $4n-3$ occur. Modulo $16$ the series $2x/(1+x^4)$ satisfies the
equation, because its fourth power vanishes and only three terms of the infinite sum survive. The results have been
checked by computer and proved formally in the Lean proof assistant.
\end{abstract}
\begin{center}\begin{minipage}{0.88\linewidth}\small
\textbf{How this paper was produced.} This paper reports the result of an autonomous AI research effort.
ChatGPT (OpenAI), using the GPT-6 model with Pro reasoning and working from prompts by Dev~D.~Goyal, found the proof, wrote its Lean~4 formalization, and wrote a first draft of the paper.
Dev~D.~Goyal chose the problems and asked for checks and revisions, but did not supply proof ideas.
Claude Opus~5.5, used with Claude Code, then checked the work independently: it rebuilt the formal proofs from
source and confirmed that they rely on no unproved assumptions, recomputed the sequence from its definition and tested
every claim against the published data, and read every step of the proofs. It also rewrote this paper for clarity.
Dev~D.~Goyal has verified the results and is responsible for them. The paper has not been peer reviewed.
The section ``About this paper'' at the end gives the details.
\end{minipage}\end{center}

\section{Introduction}

\subsection{The sequence}\label{sec:sequence}
The \emph{On-Line Encyclopedia of Integer Sequences} (OEIS)~\cite{oeisdb} is a public database of integer sequences.
Besides the terms, many entries record formulas and \emph{conjectures}: patterns observed in the data that the
contributor could not prove. This paper concerns entry A355872~\cite{oeis}, contributed by Paul~D.~Hanna in 2022, which
reads in part:
\begin{quote}\small\ttfamily\raggedright G.f. A(x) satisfies: x = Sum\_\{n=-oo..+oo\} (-x)\^{}(n\^{}2) * A(x)\^{}((n-1)\^{}2).\par\smallskip Conjecture: a(n) == 2 (mod 4) for n >= 1.\par\smallskip Conjecture: a(2*n-1) == 2 (mod 8) for n >= 1.\par\smallskip Conjecture: a(2*n) == 6 (mod 8) for n >= 1.\par\smallskip G.f. A(x) = Sum\_\{n>=1\} a(n) * x\^{}(4*n-3) satisfies: ...\par\smallskip\hfill\textrm{--- OEIS A355872, Paul D. Hanna, August 2022}\end{quote}

Here ``G.f.''\ stands for \emph{generating function}, the power series whose coefficients are the terms of the
sequence. Written with $j$ for the summation index, the defining equation is
\begin{equation}\label{eq:source}
x=\sum_{j\in\Z}(-x)^{j^2}A(x)^{(j-1)^2},\qquad A(x)=\sum_{n\ge1}a(n)\,x^{4n-3}.
\end{equation}
The sum runs over \emph{all} integers $j$, positive and negative; such a sum is called \emph{bilateral}. It has the
shape of the classical theta series $\sum_{j\in\Z}q^{j^2}$, which is why equations of this kind are called
theta-type. The first terms of the sequence are 2, 14, 434, 17662, 829314, 42293582, \dots{} (from $n=1$).
Section~\ref{sec:background} explains why the infinite sum in \eqref{eq:source} makes sense and why it determines $A$.

\subsection{The question}
The three conjectures concern the remainders of $a(n)$ on division by $4$ and by $8$. The table shows the remainders
modulo $8$ and modulo $16$ for the first $32$ terms; the indices in bold are the odd ones.
\begin{center}\small\setlength{\tabcolsep}{3.2pt}
\begin{tabular}{>{\raggedright\arraybackslash}p{6.4em}>{\centering\arraybackslash}p{1.15em}>{\centering\arraybackslash}p{1.15em}>{\centering\arraybackslash}p{1.15em}>{\centering\arraybackslash}p{1.15em}>{\centering\arraybackslash}p{1.15em}>{\centering\arraybackslash}p{1.15em}>{\centering\arraybackslash}p{1.15em}>{\centering\arraybackslash}p{1.15em}>{\centering\arraybackslash}p{1.15em}>{\centering\arraybackslash}p{1.15em}>{\centering\arraybackslash}p{1.15em}>{\centering\arraybackslash}p{1.15em}>{\centering\arraybackslash}p{1.15em}>{\centering\arraybackslash}p{1.15em}>{\centering\arraybackslash}p{1.15em}>{\centering\arraybackslash}p{1.15em}}
\toprule
$n$ & \hl{1} & 2 & \hl{3} & 4 & \hl{5} & 6 & \hl{7} & 8 & \hl{9} & 10 & \hl{11} & 12 & \hl{13} & 14 & \hl{15} & 16 \\ \midrule
$a(n)\bmod 8$ & \hl{2} & 6 & \hl{2} & 6 & \hl{2} & 6 & \hl{2} & 6 & \hl{2} & 6 & \hl{2} & 6 & \hl{2} & 6 & \hl{2} & 6 \\
$a(n)\bmod 16$ & \hl{2} & 14 & \hl{2} & 14 & \hl{2} & 14 & \hl{2} & 14 & \hl{2} & 14 & \hl{2} & 14 & \hl{2} & 14 & \hl{2} & 14 \\
\bottomrule
\end{tabular}
\\[4pt]
\begin{tabular}{>{\raggedright\arraybackslash}p{6.4em}>{\centering\arraybackslash}p{1.15em}>{\centering\arraybackslash}p{1.15em}>{\centering\arraybackslash}p{1.15em}>{\centering\arraybackslash}p{1.15em}>{\centering\arraybackslash}p{1.15em}>{\centering\arraybackslash}p{1.15em}>{\centering\arraybackslash}p{1.15em}>{\centering\arraybackslash}p{1.15em}>{\centering\arraybackslash}p{1.15em}>{\centering\arraybackslash}p{1.15em}>{\centering\arraybackslash}p{1.15em}>{\centering\arraybackslash}p{1.15em}>{\centering\arraybackslash}p{1.15em}>{\centering\arraybackslash}p{1.15em}>{\centering\arraybackslash}p{1.15em}>{\centering\arraybackslash}p{1.15em}}
\toprule
$n$ & \hl{17} & 18 & \hl{19} & 20 & \hl{21} & 22 & \hl{23} & 24 & \hl{25} & 26 & \hl{27} & 28 & \hl{29} & 30 & \hl{31} & 32 \\ \midrule
$a(n)\bmod 8$ & \hl{2} & 6 & \hl{2} & 6 & \hl{2} & 6 & \hl{2} & 6 & \hl{2} & 6 & \hl{2} & 6 & \hl{2} & 6 & \hl{2} & 6 \\
$a(n)\bmod 16$ & \hl{2} & 14 & \hl{2} & 14 & \hl{2} & 14 & \hl{2} & 14 & \hl{2} & 14 & \hl{2} & 14 & \hl{2} & 14 & \hl{2} & 14 \\
\bottomrule
\end{tabular}
\end{center}

Modulo $16$ the remainders alternate between $2$ and $14$, and this is what we prove.

\subsection{What this paper proves}
\textbf{Main result} (Theorem~\ref{thm:main}). \emph{Equation \eqref{eq:source} has exactly one solution, only the
degrees $4n-3$ occur in it, and for every $n\ge1$}
\[
a(n)\equiv2\,(-1)^{n-1}\pmod{16}.
\]
Reducing modulo $8$ and modulo $4$ gives all three conjectures of the entry (Corollary~\ref{cor:conj}). The key step is
the identity
\[
A(x)\equiv\frac{2x}{1+x^4}=2x-2x^5+2x^9-\cdots\pmod{16}.
\]

\subsection{What was known}
The entry records an asymptotic formula for $a(n)$ (due to Vaclav Kotesovec) and two product formulas obtained from the
Jacobi triple product identity. When the entry was checked on 24~October~2026, the three congruences were still labelled
as conjectures. We are not aware of an earlier proof.

\subsection{Why it is interesting}
Theta-type sums such as $\sum_{j\in\Z}q^{j^2}z^j$ are among the most studied objects in combinatorics and number theory;
through the Jacobi triple product identity they are connected with partitions and with sums of squares. Here such a sum
is used to define a power series implicitly, and the coefficients that result grow quickly (roughly like $70^n$). The
congruences show that, modulo $16$, all of this complexity collapses to a two-term periodic pattern. The proof
illustrates a mechanism that works in other settings: modulo a power of $2$, a series all of whose coefficients are
even has a power equal to zero, so an infinite sum can reduce to finitely many terms.

\subsection{How to read this paper}
Section~\ref{sec:background} recalls power series, congruences and bilateral sums. Section~\ref{sec:results} states the
results and Section~\ref{sec:proofs} proves them. Section~\ref{sec:checks} describes the computer checks and the formal
proof, and Section~\ref{sec:discussion} discusses the scope of the result.
\section{Background}\label{sec:background}
This section collects the definitions and facts used later. Readers who know formal power series and congruences can skip it and return when needed.
\subsection{Formal power series}
A \emph{formal power series} is an expression
\[
U(x)=u_0+u_1x+u_2x^2+u_3x^3+\cdots ,
\]
determined by its sequence of coefficients $u_0,u_1,u_2,\dots$. No number is ever substituted for $x$, so questions
of convergence do not arise. We write $\coeff{n}{U}=u_n$ for the coefficient of $x^n$ in $U$; two series are equal
when all their coefficients are equal.

Series are added coefficient by coefficient and multiplied like polynomials:
\[
\coeff{n}{(UV)}=u_0v_n+u_1v_{n-1}+\cdots+u_nv_0 .
\]
Each coefficient of a product is a finite sum. For example, $(1-x)(1+x+x^2+x^3+\cdots)=1$, so $1+x+x^2+\cdots$ is
the \emph{multiplicative inverse} of $1-x$, written $\frac{1}{1-x}$. More generally, every series $D$ with constant term
$1$ has a unique multiplicative inverse $\frac1D$, whose coefficients are determined one at a time by $D\cdot\frac1D=1$
using only addition, subtraction and multiplication. Every fraction in this paper has a denominator of this kind, so it
denotes a power series.

We write $U=O(x^d)$ when all coefficients of $U$ below degree $d$ are zero. So
$U-V=O(x^d)$ means that $U$ and $V$ agree up to degree $d-1$. A series of the form $x+u_2x^2+u_3x^3+\cdots$, with
constant term $0$ and coefficient of $x$ equal to $1$, is called \emph{normalized}.

\subsection{Remainders and congruences}
For an integer $m\ge2$ we write $a\equiv b\pmod m$ (``$a$ is congruent to $b$ modulo $m$'') when $a-b$ is divisible by
$m$, that is, when $a$ and $b$ leave the same remainder after division by $m$. For example, $17\equiv5\pmod 4$, since both leave remainder $1$, and $-3\equiv1\pmod4$. The
remainder of $a$ on division by $m$ is called its \emph{residue} modulo $m$; it is one of $0,1,\dots,m-1$. Saying that
$a$ is odd is the same as saying $a\equiv1\pmod2$.

Addition and multiplication are compatible with congruences, so the residues $0,1,\dots,m-1$, added and multiplied
modulo $m$, form a ring, written $\Z/m\Z$. When $m=p$ is a prime we also write $\F_p$. In $\F_2=\{0,1\}$ we have $1+1=0$,
so addition and subtraction coincide.

For power series with integer coefficients, $U\equiv V\pmod m$ means that \emph{every} coefficient of $U-V$ is
divisible by $m$. Equivalently, the two series become equal when every coefficient is replaced by its residue; we call
this \emph{reducing} the series modulo $m$. The key fact is the following.

\emph{Reduction modulo $m$ is compatible with all the operations above.} To reduce a sum, a product, a substitution, or
an inverse, we may reduce the ingredients first and then compute in $\Z/m\Z$. This holds because every
coefficient of the result is obtained from finitely many coefficients of the ingredients by adding, subtracting and
multiplying, and these operations are compatible with remainders.

\para{Rings.} Several lemmas below hold over every \emph{commutative ring}, that is, every set with an addition,
subtraction and multiplication satisfying the usual rules of arithmetic (for example $\Z$, $\Q$ and $\Z/m\Z$). No
division is used; this matters because in $\Z/m\Z$ division is not always possible (for instance, $2$ has no inverse in
$\Z/4\Z$).

\subsection{Sums over all integers}
For a series $F$ with coefficients in a commutative ring $R$ and an integer $j$, put
\[
T_j(F)=(-x)^{j^2}F^{(j-1)^2}.
\]
Both exponents are nonnegative integers, so $T_j(F)$ is a power series. Because of the factor $(-x)^{j^2}$, its
coefficients of degree less than $j^2$ are zero. Hence a given coefficient of degree $d$ receives contributions only
from the finitely many $j$ with $j^2\le d$, and the bilateral sum
\[
\Theta(F)=\sum_{j\in\Z}T_j(F)
\]
is well defined coefficient by coefficient. Equation \eqref{eq:source} says that $\Theta(A)=x$. No question of
convergence arises: every coefficient of $\Theta(F)$ is a finite sum.

The entry also gives the equation with $(j+1)^2$ in place of $(j-1)^2$; replacing $j$ by $-j$ shows that the two forms
are the same.

\section{Results}\label{sec:results}

\begin{theorem}\label{thm:main}
\begin{enumerate}[label=(\alph*)]
\item Over every commutative ring $R$ there is exactly one power series $A_R$ with $\Theta(A_R)=x$.
\item The integer solution $A=A_\Z$ has nonzero coefficients only in degrees $4n-3$, $n\ge1$; so
$A(x)=\sum_{n\ge1}a(n)x^{4n-3}$ as in \eqref{eq:source}.
\item Modulo $16$,
\[
A(x)\equiv\frac{2x}{1+x^4}\pmod{16},\qquad\text{equivalently}\qquad a(n)\equiv2\,(-1)^{n-1}\pmod{16}\quad(n\ge1).
\]
\end{enumerate}
\end{theorem}

\begin{corollary}[Hanna's conjectures]\label{cor:conj}
For every $n\ge1$: $a(n)\equiv2\pmod4$, $a(2n-1)\equiv2\pmod8$, and $a(2n)\equiv6\pmod8$.
\end{corollary}

\begin{proof}
By Theorem~\ref{thm:main}(c), $a(n)$ leaves remainder $2$ on division by $16$ when $n$ is odd and remainder $14$ when $n$
is even. Reducing, the remainder modulo $8$ is $2$ for odd $n$ and $6$ for even $n$, and modulo $4$ it is $2$ in both
cases.
\end{proof}

\section{Proofs}\label{sec:proofs}

\para{The idea in brief.} The term with $j=0$ in the bilateral sum is $A$ itself. Moving all other terms to the other
side expresses $A$ through terms that each carry at least one extra factor of $x$, so the equation determines the
coefficients one at a time, over every ring. Modulo $16$, the series $P=2x/(1+x^4)$ has all coefficients even, so
$P^4=0$; in $\Theta(P)$ only the terms with $(j-1)^2<4$, that is $j=0,1,2$, survive, and they add up to $x$.

Throughout, $R$ is a commutative ring and $\Phi(F)=x-\bigl(\Theta(F)-F\bigr)$, so that $\Theta(F)=x$ is equivalent to
$F=\Phi(F)$.

\begin{lemma}\label{lem:contract}
Let $r\ge0$ and $F,G\in\ps{R}$. If $F-G=O(x^r)$, then $\Phi(F)-\Phi(G)=O(x^{r+1})$.
\end{lemma}

\begin{proof}
For every $e\ge0$, $F^e-G^e=(F-G)(F^{e-1}+F^{e-2}G+\cdots+G^{e-1})$ is a multiple of $F-G$, hence $O(x^r)$. For
$j\neq0$ we have $j^2\ge1$, so $T_j(F)-T_j(G)=(-x)^{j^2}\bigl(F^{(j-1)^2}-G^{(j-1)^2}\bigr)=O(x^{r+1})$. The term
$j=0$, which equals $F-G$, does not occur in $\Phi(F)-\Phi(G)$. Each coefficient of $\Phi(F)-\Phi(G)$ is a finite sum of
coefficients of the differences $T_j(F)-T_j(G)$ with $j\ne0$, which proves the claim.
\end{proof}

\begin{proof}[Proof of Theorem~\ref{thm:main}(a)]
By Lemma~\ref{lem:contract}, the coefficient of $x^d$ in $\Phi(F)$ depends only on the coefficients of $F$ of degree
below $d$. Choosing the coefficients of $A_R$ one at a time so that $A_R=\Phi(A_R)$ holds in each degree therefore gives
a solution, and two solutions must agree in every degree. Equivalently, starting from $F_0=0$ and putting
$F_{k+1}=\Phi(F_k)$, Lemma~\ref{lem:contract} gives $A_R-F_k=O(x^k)$, so each coefficient of $F_k$ is constant from
some $k$ on and equals the corresponding coefficient of $A_R$.
\end{proof}

No division was used, so reducing the integer solution modulo $m$ gives a solution over $\Z/m\Z$, and by uniqueness it
is the solution there (Section~\ref{sec:background}).

\begin{proof}[Proof of Theorem~\ref{thm:main}(b)]
Say that a series is \emph{supported in class $1$} if its only nonzero coefficients are in degrees $\equiv1\pmod4$. If
$F$ is supported in class $1$, then $F^e$ is supported in degrees $\equiv e\pmod4$, and $T_j(F)$ in degrees
$\equiv j^2+(j-1)^2\pmod4$. Of the two consecutive integers $j-1$ and $j$, one is even and one is odd, so
$j^2+(j-1)^2\equiv0+1=1\pmod4$. Hence $\Theta(F)$, and therefore $\Phi(F)$, is again supported in class $1$. All the
approximations $F_k$ from the proof of part (a) are thus supported in class $1$, and so is their limit $A$. This also
shows that $A$ has constant term $0$.
\end{proof}

\begin{proof}[Proof of Theorem~\ref{thm:main}(c)]
Work over $R=\Z/16\Z$ and put
\[
P(x)=\frac{2x}{1+x^4}=2\sum_{k\ge0}(-1)^kx^{4k+1}.
\]
Every coefficient of $P$ is even, so every coefficient of $P^4$ is divisible by $2^4=16$; that is, $P^4=0$ in $\ps{R}$,
and then $P^e=0$ for all $e\ge4$. Now $(j-1)^2\ge4$ unless $j\in\{0,1,2\}$, so all terms of $\Theta(P)$ vanish except
\[
T_0(P)=P,\qquad T_1(P)=-x,\qquad T_2(P)=(-x)^4P=x^4P .
\]
Hence
\[
\Theta(P)=P-x+x^4P=(1+x^4)P-x=2x-x=x .
\]
By part (a), $P$ is the unique solution over $\Z/16\Z$, so it is the reduction of $A$ modulo $16$. The coefficient of
$x^{4n-3}=x^{4(n-1)+1}$ in $P$ is $2(-1)^{n-1}$, which gives the formula for $a(n)$.
\end{proof}

\section{Checking the results}\label{sec:checks}

A mathematical proof can contain mistakes, so the results were checked in two independent ways.

\para{Numerical check.}
We computed the solution of \eqref{eq:source} by the fixed-point iteration of the proof, exactly up to degree $120$
(that is, $n\le30$) and modulo $16$ up to degree $1200$. The computation reproduces all 17 terms printed in the OEIS
entry, confirms that only the degrees $4n-3$ occur, and confirms $a(n)\equiv2(-1)^{n-1}\pmod{16}$ for $n\le300$.
The computation was done with the PARI/GP computer algebra system, using code that was written separately from the
formal proofs (file \texttt{verification/independent\_checks.gp} in the accompanying repository~\cite{repo}).
A computation can only test finitely many cases, so it is a check on the definitions and on the indexing, not a proof.

\para{Formal proof in Lean.}
Lean~\cite{lean} is a \emph{proof assistant}: a program in which mathematical statements and their proofs are written
in a precise formal language, and every step is checked mechanically down to the basic rules of logic. All results of
this paper are proved in Lean~4 (version 4.19.0) on top of Mathlib~\cite{mathlib}, Lean's library of formalized
mathematics. The statements are collected in the file \texttt{lean/Entries/A355872.lean} of the repository,
in the namespace \texttt{OEIS.A355872}. The construction over an arbitrary commutative ring, its compatibility with reduction of coefficients,
and the congruence for the whole series modulo $16$ are in \texttt{lean/Families/BilateralTheta.lean} (namespace
\texttt{Families.BilateralTheta}). The hypotheses of the theorems in the table are the literal equation
\eqref{eq:source}; no residue is assumed.
The table lists where each result of this paper is proved.
\begin{center}\small\renewcommand{\arraystretch}{1.25}
\begin{tabular}{>{\raggedright\arraybackslash}p{0.52\linewidth}>{\raggedright\arraybackslash}p{0.4\linewidth}}
\toprule
Statement in this paper & Lean theorem(s)\\
\midrule
Equation~\eqref{eq:source}, written coefficient by coefficient, is equivalent to $\Theta(A)=x$ & \texttt{literal\_\allowbreak{}source\_\allowbreak{}iff} \\
Theorem~\ref{thm:main}(a) over $\Z$ & \texttt{source\_\allowbreak{}existence\_\allowbreak{}uniqueness}, \texttt{canonical\_\allowbreak{}literal\_\allowbreak{}source} \\
Theorem~\ref{thm:main}(b) & \texttt{canonical\_\allowbreak{}sparse\_\allowbreak{}support} \\
Theorem~\ref{thm:main}(c), for every integer series satisfying \eqref{eq:source} & \texttt{source\_\allowbreak{}mod\_\allowbreak{}sixteen} \\
Corollary~\ref{cor:conj} & \texttt{source\_\allowbreak{}all\_\allowbreak{}three} \\
\bottomrule
\end{tabular}
\end{center}
The Lean project was rebuilt from source for this paper, and Lean was asked to list the basic assumptions (axioms)
that each of these theorems ultimately depends on. Every one depends only on the three standard axioms of Lean's logic
(\texttt{propext}, \texttt{Classical.choice}, \texttt{Quot.sound}); no step is left unproved. A formal proof certifies
exactly the statement that Lean checks. Whether that statement says what the OEIS entry means is a separate question,
so the Lean definitions were also compared, side by side, with the definition in the OEIS entry
(Section~\ref{sec:sequence}).

\section{Discussion}\label{sec:discussion}

\para{Scope.} Theorem~\ref{thm:main} determines every coefficient modulo $16$, so it settles the three conjectures of the
entry completely. It says nothing about the size of $a(n)$; the asymptotic formula and the product formulas recorded in
the entry are not addressed.

\para{Higher powers of $2$.} The same mechanism does not immediately give more: modulo $32$ the fourth power of $P$ is
no longer zero, and more terms of the bilateral sum survive. A computer calculation suggests that the residues remain
periodic with period $4$: for $n<200$, $a(n)$ modulo $32$ repeats $2,14,18,30$, and modulo $64$ it repeats
$2,14,50,62$, both from $n=1$. It would be interesting to describe $a(n)$ modulo every power of $2$.

\para{Related equations.} The argument of this paper applies to any bilateral equation in which, as here, the term
with $j=0$ is the unknown series and all other terms carry a positive power of $x$: existence and uniqueness then hold
over every ring, and congruences can be found by exhibiting a solution of the reduced equation.

\section*{About this paper}\label{sec:about}
\para{How the work was done.}
This paper is one of 29 papers on conjectures from the OEIS that came out of an autonomous AI research effort directed
by Dev~D.~Goyal~\cite{repo}. The roles were as follows.
\begin{itemize}
\item \emph{Dev~D.~Goyal} started the project, chose which problems to pursue, asked for proofs, formal verification and
revisions, and decided what to publish. Dev~D.~Goyal did not supply proof strategies or intermediate arguments.
\item \emph{ChatGPT} (OpenAI), using the GPT-6 model with Pro reasoning, worked autonomously from those prompts. It
searched the OEIS for suitable conjectures, found the proofs, wrote the Lean~4 formalization, and wrote a first draft
of this paper.
\item \emph{Claude Opus~5.5}, used with Claude Code, checked the work independently and then rewrote this paper. It
rebuilt the Lean project from source and confirmed that the theorems listed in Section~\ref{sec:checks} rely only on
Lean's standard axioms; compared the Lean definitions with the OEIS definition; recomputed the sequence with separately
written PARI/GP code, compared it with the terms published in the OEIS, and tested every claim of the paper on many
terms; and read every step of the proofs. It then rewrote the text to make it self-contained and readable. The
theorems and the proof strategies are ChatGPT's. In rewriting, Claude Opus~5.5 expanded the explanations and, in
places, wrote out an argument in more detail or by a more elementary route. The Discussion section was added by Claude Opus~5.5;
the further observations in it were checked by computer and are not part of the Lean formalization.
\end{itemize}
Dev~D.~Goyal has verified the results and is responsible for this paper.

\para{Review status.} This paper has not been peer reviewed, and it has not been submitted to a journal or to the OEIS.

\para{Priority.} An OEIS entry that labels a statement as a conjecture shows only that its contributor did not know
a proof; it does not show that no proof existed. During the research only limited searches for earlier work were
made. No claim of priority is made.

\para{Corrections.} Please report any error in the issue tracker of the repository~\cite{repo}, where the Lean sources and the numerical checks are also
available.

\section*{Acknowledgements}
The sequence and its conjectures are due to Paul~D.~Hanna~\cite{oeis}. The asymptotic formula mentioned in the
introduction was added to the entry by Vaclav~Kotesovec.

\begin{thebibliography}{99}\small
\raggedright
\bibitem{oeis} P.~D.~Hanna, entry A355872 (G.f.\ $A(x)$ satisfies: $x=\sum_{n\in\Z}(-x)^{n^2}A(x)^{(n-1)^2}$), \emph{The On-Line Encyclopedia of Integer Sequences}, August 2022. \url{https://oeis.org/A355872}. Text as of October 2026.
\bibitem{oeisdb} N.~J.~A.~Sloane and The OEIS Foundation Inc., \emph{The On-Line Encyclopedia of Integer Sequences}. \url{https://oeis.org}.
\bibitem{stanley} R.~P.~Stanley, \emph{Enumerative Combinatorics}, Volume~2, Cambridge University Press, 1999. (Chapter~5 treats formal power series, composition and compositional inverses.)
\bibitem{niven} I.~Niven, Formal power series, \emph{American Mathematical Monthly} 76 (1969), 871--889.
\bibitem{lean} L.~de~Moura and S.~Ullrich, The Lean~4 theorem prover and programming language, in \emph{Automated Deduction -- CADE 28}, Lecture Notes in Computer Science 12699, Springer, 2021, 625--635.
\bibitem{mathlib} The mathlib Community, The Lean mathematical library, in \emph{Proceedings of the 9th ACM SIGPLAN International Conference on Certified Programs and Proofs (CPP 2020)}, 367--381. Version used: commit \texttt{c44e0c8ee63c}.
\bibitem{repo} D.~D.~Goyal, \emph{Proofs and corrections of OEIS conjectures, with Lean~4 formalizations}, repository, 2026. \url{https://github.com/dvgyl/oeis-a396843-proof}.
\end{thebibliography}
\end{document}
